In dynamic programming, define what information a state retains before writing transitions and boundary values. Limits on deletions, inventory, or selected intervals must be represented in the state. Unreachable states must be distinguished from feasible states with zero reward. In counting problems, different transitions must describe disjoint sets of solutions to avoid double counting.

The choice of algorithm rests on a proved structure: binary search needs a monotone exclusion rule, greedy methods need an exchange or dominance argument, and divide-and-conquer needs a way to account for contributions across subproblems. Count all states, transitions, and data-structure operations when bounding complexity. Space compression also requires an update order that does not accidentally reuse current-round values. For randomized algorithms, distinguish distributional correctness, expected cost, and worst-case cost.
\par\smallskip
